\section{Discussion}

The cutting voxel projector (CVP) algorithm, based on the principles in \cref{sec:cuts}, was first implemented in 2019. Since its inception, it has been employed in hundreds of reconstruction problems, particularly in improving algorithms for processing CBCT perfusion data. Notably, it has contributed to prior works such as \cite{Kulvait2021, Haseljic2021, Kulvait2022}. Throughout its applications, the projector has consistently demonstrated high accuracy, comparable to or exceeding the established TT footprint projector, especially for complex cone-beam geometries. 

More recently, the algorithm has been adapted for parallel-ray geometry, where the absence of elevation correction requirements makes it an almost exact projector, e.g. for synchrotron $\mu$-CT setups, see also \cite{Moosmann2014}. This extension underscores its versatility and accuracy across a broader range of tomographic applications. The GitHub repository with the implementation is often cloned and has 19 stars as of November 2024.

\subsection{Elevation correction}
When comparing projector profiles across elevation angles, as seen in \cref{fig:projectorerror}, the CVP exhibited consistently low errors but occasionally showed localized spikes above the TT projector. Closer examination revealed that these errors occurred at the top or bottom edges of voxel projections and were associated with how the center of mass of the voxel's slice aligned with pixel boundaries in the $\chi_2$ direction.

To address this, we introduced an elevation correction. This modification extends the projected line (defined by the center of mass and $x_3$ limits) into a rectangle, allowing a portion of the projection to be assigned to the appropriate pixel. While this increases computational complexity, it also improves accuracy without a substantial impact on runtime, particularly for the relaxed version of the projector.

For most use cases, especially those with small cone angles, the CVP without elevation correction remains a very good choice, offering high accuracy with fast performance. However, for applications with large cone angles or precise scenarios requiring maximum accuracy, the elevation correction provides clear benefits, as shown in \cref{fig:projectorerror}.

\subsection{Speed and GPU optimization}
Tables \ref{speed:long} and \ref{speed:pfeiffertest} highlight the impressive runtime performance of the CVP, particularly its relaxed variant. Achieving these results required GPU-specific optimizations, including computing cuts in the $x_1x_2$ plane separately for each voxel. This strategy aligns well with GPU architectures, enabling effective parallelization by assigning one thread per voxel. Furthermore, optimized memory conflict handling in the local GPU memory was implemented to minimize memory access contention and maximize computational throughput.

Despite these optimizations, the projector remains slower than the backprojector. For speed sensitive use cases, it may be advisable to relax the accuracy requirements of the forward projector and use a fast, low-accuracy ray-tracing method. The backprojector, however, cannot be similarly simplified, as doing so would amplify operator inaccuracies, particularly in iterative reconstruction methods. Frameworks like ASTRA, see \cite{ASTRA2016}, often adopt this strategy, combining a forward raytracer and a backward footprint projector for speed efficiency, even at the cost of sacrificing the self-adjointness of the tomographic operator.

\subsection{Possible Improvements}
While the CVP algorithm has been extensively optimized, these improvements are possible in case of footprint projectors as well. For example, applying the same parallelization techniques used in the CVP to footprint projectors, such as the TT or TR projectors, could significantly accelerate their performance. Current implementations, such as the approach described by Long et al. \cite{Long2010}, are primarily CPU-optimized and do not fully exploit GPU capabilities. By restructuring the calculations to parallelize over voxels rather than rows and columns, as demonstrated in our CVP implementation, significant speedups could likely be achieved for these projectors as well.

%Additionally, for applications requiring extremely precise reconstructions and involving large cone angles, the role of elevation correction becomes increasingly important. Future work will investigate the performance of the elevation-corrected CVP in such scenarios to determine its potential benefits in these high-accuracy use cases.

In future work, we plan to explore the application of the CVP to non-circular trajectories. While the current implementation in \cref{sec:implementation_full} assumes alignment between the voxel $x_3$ axis and the detector $\chi_1$ axis, the derivation in \cref{sec:formulation} is independent of this. By leveraging algorithms for computing general cube and square pyramid intersections, it would be possible to construct projectors for arbitrary scanning trajectories, distinguishing the CVP from footprint projectors that inherently rely on this geometric alignment by design.

\subsection{Conclusion}
This work introduces and evaluates a cutting voxel projector that balances speed and accuracy, outperforming existing projectors like TT and Siddon's algorithm for CBCT geometries. The elevation-corrected CVP is presented as the most accurate known projector, except for extremely computationally expensive ray-tracing approaches. Its GPU-optimized implementation and superior runtime performance make it a promising tool for advancing algebraic reconstruction techniques.

A key advantage of the proposed cutting voxel projector lies in its support for hierarchical reconstruction with nonuniform voxel grids. This flexibility opens the door to developing advanced computational techniques analogous to algebraic multigrid (AMG) methods used in solving partial differential equations. By enabling local voxel refinement, the projector provides the groundwork for adaptive tomographic reconstruction approaches that can focus computational resources on regions of interest while maintaining coarser resolution elsewhere. In future work, we aim to leverage this capability to design hierarchical solvers for the tomographic reconstruction problem, incorporating multigrid-inspired strategies such as dynamic grid coarsening and refinement. These developments hold the promise of significantly improving the scalability and efficiency of algebraic reconstruction techniques, particularly for large-scale or high-resolution imaging tasks. The cutting voxel projector thus not only delivers improved accuracy and performance today but also establishes a foundation for transformative innovations in adaptive tomographic reconstruction techniques.